% Compile with: uplatex file.tex && dvipdfmx file.dvi
%           or: latexmk -uplatex file.tex
%\documentclass[uplatex,dvipdfmx]{jsarticle}
\documentclass[dvipdfmx,a4paper,11pt]{jsarticle}
\usepackage{amsmath, amssymb}
\usepackage{tikz}

\title{scenario\_poisson\_bracket}
\date{\today}

\begin{document}
\maketitle

ポアソン括弧を計算するには、どの座標とどの運動量が共役（正準対）かをあらかじめ宣言しておく必要がある。動径座標 $r$ の共役運動量 $p_r$ を宣言する。

$p_r$ を r の共役運動量とする

正準座標対 (r, p\_r) を登録しました。以後「f と g のポアソン括弧を求めて」で利用できます。

角度座標 $\theta$ の共役運動量 $p_\theta$ を宣言する。

$p_{\theta}$ を theta の共役運動量とする

正準座標対 (theta, p\_theta) を登録しました。以後「f と g のポアソン括弧を求めて」で利用できます。

基本的なポアソン括弧 $\{r,p_r\}$ を求める。正準対どうしなので正準交換関係 $\{q,p\}=1$ により1になるはずである。

r と $p_r$ のポアソン括弧を求めて

\begin{equation}
  1
  \label{eq:E1}
\end{equation}

中心力ハミルトニアン $H=p_r^2/(2m)+p_\theta^2/(2mr^2)+U(r)$（動径・角度方向の運動エネルギーとポテンシャル $U(r)$ の和）を定義する。

$H = p_r^2/(2 m) + p_{\theta}^2/(2 m r^2) + U(r)$

\begin{equation}
  H = p_r^2/(2 m) + p_{\theta}^2/(2 m r^2) + U(r)
  \label{eq:E2}
\end{equation}

角運動量 $L_z=p_\theta$（角度の共役運動量そのもの）を定義する。

$L_z = p_{\theta}$

\begin{equation}
  L_z = p_{\theta}
  \label{eq:E3}
\end{equation}

ハミルトニアン $H$ と角運動量 $L_z$ のポアソン括弧を求める。$H$ が角度 $\theta$ に依らない（回転対称）ので、これがゼロになれば $L_z$ は保存量であることが確認できる。

H と $L_z$ のポアソン括弧を求めて

\begin{equation}
  0
  \label{eq:E4}
\end{equation}

ハミルトンの運動方程式 $\dot q=\{q,H\}$ を用いて動径座標 $r$ の時間発展を求めるため、$H$ と $r$ のポアソン括弧を求める。

H と r のポアソン括弧を求めて

\begin{equation}
  -\frac{p_{r}}{m}
  \label{eq:E5}
\end{equation}

\end{document}